% Generated by task tex-groundtruth from dossiers/gt-bus.md (section A, A.1 to A.5),
% with the corrections of dossiers/verify-gt-bus.md.
\section{The setup, the commitment and the bus phase: the transcript as deployed}\label{sec:gen-bus-transcript}

\begin{sloppypar}
Short forms in this section. Specification files are under \file{doc/leanvm/body/}:
\code{03:} is \file{03-proving-primitives.tex}, \code{05:} is \file{05-arithmetization.tex},
\code{06:} is \file{06-bus-interactions.tex}, \code{08:} is \file{08-end-to-end-protocol.tex}.
Rust files are under \file{crates/lean_vm/src/}: \code{mod.rs} is \file{cpu/mod.rs},
\code{layout.rs} is \file{cpu/layout.rs}, and \code{leaf.rs}, \code{gkr.rs}, \code{witness.rs},
\code{tables.rs}, \code{pcs.rs}, \code{constraints.rs} are at that root; \code{fs/} is
\file{crates/fiat_shamir/src/}, and \code{primitives/src/} is \file{crates/primitives/src/}.
\code{py:} is \file{python-verifier/verifier.py}. The number after the colon is the line, at
leanVM's pin \code{a386121f}.
\end{sloppypar}

\subsection{Notation and sizes}\label{sec:gen-bus-notation}

\begin{itemize}
\item $\K = \mathrm{GF}(2^{64})$, $\E = \mathrm{GF}(2^{192})$; every stream element is one element
  of $\E$ (24 bytes).
\item Announced: $κ_{\mathrm{mem}}$ (memory log-size), $τ_0 … τ_5$ (log-heights of \code{XOR},
  \code{MUL}, \code{SET}, \code{DEREF}, \code{JUMP}, \code{BLAKE2S}), $ρ$ (the rate exponent
  \code{log_inv_rate}). Public: the program, of $2^{κ_{\mathrm{bc}}}$ instructions, and the public
  input, two words.
\item $μ_{\mathrm{stack}}$: log-size of the committed stack. $μ$ (written $μ_{\mathrm{bus}}$ where
  confusion is possible): depth of the three leaf trees of the bus,
  $μ = ⌈\log_2 N_{\mathrm{push}}⌉$ with
  $N_{\mathrm{push}} = 1 + 2^{κ_{\mathrm{mem}}} + 2^{κ_{\mathrm{bc}}} + Σ_j \mathit{nfl}_j · 2^{τ_j}$,
  where $\mathit{nfl} = (5, 5, 3, 5, 5, 11)$ is the number of flush pairs of each table
  (\code{tables.rs:483-495, 548-561, 633-643, 731-751, 873-908}). For leanVM $17 ≤ μ ≤ 28$ (the
  lower bound from $κ_{\mathrm{mem}} ≥ 16$; the upper from $μ_{\mathrm{stack}} ≤ 28$, dossier
  \file{gt-bus.md}, section D.4).
\item The count side has $\mathit{ncnt} = (4, 4, 2, 4, 4, 10)$ blocks per table, 28 in all
  (\code{tables.rs:479-482, 544-547, 629-632, 718-721, 869-872}).
\end{itemize}

\subsection{Setup and commitment}\label{sec:gen-bus-setup}

$P→V$ is a prover message (a read of the stream, bound into the Fiat--Shamir state as it is
read, \code{fs/transcript.rs:283-287}), $V$ a verifier computation or check, $V→P$ a challenge
(one squeeze of the chain per element of $\E$, \code{fs/lib.rs:95-99}).

{\footnotesize\setlength{\tabcolsep}{4pt}
\begin{longtable}{@{}>{\raggedright\arraybackslash}p{0.04\linewidth}>{\raggedright\arraybackslash}p{0.28\linewidth}>{\raggedright\arraybackslash}p{0.21\linewidth}>{\raggedright\arraybackslash}p{0.24\linewidth}>{\raggedright\arraybackslash}p{0.15\linewidth}@{}}
\caption{Setup and commitment, as deployed.}\label{tab:gen-bus-setup}\\
\toprule
\# & Step & Specification & Rust & Python \\
\midrule
\endfirsthead
\toprule
\# & Step & Specification & Rust & Python \\
\midrule
\endhead
\bottomrule
\endfoot
S0 & $V$: seed the chain with \code{compress(iv, public input)},
  \code{iv = BLAKE2s("leanvm" ‖ len ‖ R1CS_DIGEST ‖ bytecode_hash)}
  & \code{08:55} ``seeded by the public input, the bytecode and Flock's BLAKE2s R1CS matrices'';
  §8.4 is \code{TODO} (\code{08:44-46})
  & \code{mod.rs:82-93}, \code{mod.rs:712}; \code{fs/lib.rs:69-73}
  & \code{py:1366-1369}, \code{py:347-351} \\
S1 & $P→V$: 8 elements: $κ_{\mathrm{mem}}$, $τ_0 … τ_5$, $ρ$, each a small integer in the low limb
  & \code{08:55} ``The prover sends the memory log-size $κ_{\mathrm{mem}}$ and the six table
  log-heights $τ_j$'' (no $ρ$)
  & prover \code{mod.rs:118-124}; verifier \code{mod.rs:144-149}
  & \code{py:1372-1376} \\
S2 & $V$: each announced element has its two upper limbs zero
  & absent
  & \code{mod.rs:133-135} \code{if word.c1 != 0 || word.c2 != 0 { return Err(CpuError::PublicInput) }}
  & \code{py:1373} \\
S3 & $V$: each public word has its third limb zero
  & \code{08:29} ``two 192-bit words (with top limb 0)''
  & \code{mod.rs:141-143}
  & by type: the input is a 32-byte \code{Digest} (\code{py:204-219}) \\
S4 & $V$: the caps (the checks B4 to B9 of \cref{tab:gen-bus-checks})
  & \code{06:80}, \code{06:86}, \code{06:90}
  & \code{mod.rs:157-170}
  & \code{py:857-864}, \code{py:1377} \\
S5 & $V$: derive the layouts (the stack, the three leaf stacks)
  & \code{08:55} ``derives the layout of every stack (PCS and GKR leaves)''
  & \code{mod.rs:171}, \code{layout.rs:328-427}, \code{witness.rs:67-102}
  & \code{py:856-895}, \code{py:506-511} \\
S6 & $V$: the stacking window $15 ≤ μ_{\mathrm{stack}} ≤ 28$
  & absent
  & \code{mod.rs:174-176}; \code{pcs.rs:49, 51}
  & \code{py:1379}; \code{py:907-908} \\
C1 & $P→V$: the commitment, a Merkle root as 2 elements (two 128-bit halves)
  & \code{08:60-61} ``sends its commitment''
  & prover \code{pcs.rs:118-119}, \code{mod.rs:560-562}; verifier \code{mod.rs:714},
  \code{pcs.rs:136-138}, \code{fs/transcript.rs:97-99}
  & \code{py:1382} \\
C2 & $V$: both halves of the root have a zero third limb
  & absent
  & \code{fs/merkle.rs:26-29} (\code{NonCanonicalEncoding})
  & \code{py:223} \\
\end{longtable}
}

No challenge is drawn before the root is read (\code{pcs.rs:93-96} ``bind its root into the
transcript, before any challenge is sampled''; the first \code{sample} is \code{leaf.rs:876}).

\subsection{The bus phase}\label{sec:gen-bus-phase}

{\footnotesize\setlength{\tabcolsep}{4pt}
\begin{longtable}{@{}>{\raggedright\arraybackslash}p{0.04\linewidth}>{\raggedright\arraybackslash}p{0.30\linewidth}>{\raggedright\arraybackslash}p{0.18\linewidth}>{\raggedright\arraybackslash}p{0.24\linewidth}>{\raggedright\arraybackslash}p{0.16\linewidth}@{}}
\caption{The bus phase, as deployed.}\label{tab:gen-bus-phase}\\
\toprule
\# & Step & Specification & Rust & Python \\
\midrule
\endfirsthead
\toprule
\# & Step & Specification & Rust & Python \\
\midrule
\endhead
\bottomrule
\endfoot
U0 & $V$: structural assertions on the leaf layouts (panics, not rejections)
  & absent
  & \code{leaf.rs:875}, \code{leaf.rs:123-146}
  & absent \\
U1 & $V→P$: $α_0, α_1, α_2, α_3$ then $β$: five squeezes, no message between them
  & \code{08:67}; \code{05:22}, \code{05:28}
  & \code{leaf.rs:876}, \code{leaf.rs:882}
  & \code{py:542}, \code{py:544} \\
U2 & $P→V$: 2 elements, the bus root $R$ then the count root $R_c$
  & \code{08:68} ``Sends the count root $R_c$ and one bus root $R$'' (prose order reversed)
  & prover \code{gkr.rs:272-276}; verifier \code{gkr.rs:364-367}
  & \code{py:434}, \code{py:437} \\
U3 & $V→P$: the first combiner $λ$
  & \code{05:91}
  & \code{gkr.rs:369}
  & \code{py:435} \\
U4 & The batched GKR (\cref{sec:gen-bus-gkr}): messages, challenges and one check per layer
  & \code{05:61-95}, \code{08:69}
  & \code{gkr.rs:373-427}
  & \code{py:439-455} \\
U5 & $V$: $R_c ≠ 0$ (evaluated after the GKR)
  & \code{06:78}, \code{08:69}, \code{08:100}
  & \code{leaf.rs:884-889} (\code{Error::ZeroCount})
  & \code{py:546} \\
U6 & $P→V$: 5 elements, the evaluations of \code{mem_0}, \code{mem_1}, \code{mem_2},
  \code{cntfin_mem} at $ζ_{<κ_{\mathrm{mem}}}$ and of \code{cntfin_bc} at $ζ_{<κ_{\mathrm{bc}}}$,
  in that order
  & \code{08:70} ``For their committed columns the prover sends one evaluation each'';
  \code{05:111}
  & verifier \code{leaf.rs:909-920}, \code{leaf.rs:428-439}, \code{leaf.rs:548-550}; prover
  \code{leaf.rs:500-529}
  & \code{py:552-557} \\
U7 & $V$: computes, per side, the public part of the leaf claim (selectors, constants, index
  column, bytecode column, the five evaluations, the padding) and \textbf{derives} what the tables
  owe, $\mathit{rem}_s$; builds the tables' bus forms. No check.
  & \code{08:70}; \code{05:105-123}
  & \code{leaf.rs:389-461}, \code{leaf.rs:921-924}
  & \code{py:559-590} \\
\end{longtable}
}

After the last of these steps (U7) the table sumcheck starts with the challenge $ξ$
(\code{mod.rs:726}, \code{py:1389}). Nothing of the bus phase is sent after the boundary
evaluations (U6).

\paragraph{Element counts} (checked by the dossier's probe, dossier \file{gt-bus.md}, section
H.1, against the pinned Python function, for $μ = 0 … 12$; the citation check re-ran the probe,
recounted every message from the code, and agrees):

\begin{center}\small
\begin{tabular}{@{}ll@{}}
\toprule
Message & Elements of $\E$ \\
\midrule
announced sizes & 8 \\
commitment root & 2 \\
roots & 2 \\
GKR & $μ² + 4μ$, plus 1 if $μ$ is odd \\
boundary evaluations & 5 \\
\textbf{total, seed to end of bus} & $μ² + 4μ + 17$ ($+1$ if $μ$ is odd) \\
\bottomrule
\end{tabular}
\qquad
\begin{tabular}{@{}ll@{}}
\toprule
Challenge & Squeezes \\
\midrule
$α$, $β$ & 5 \\
GKR, $μ$ even & $μ²/4 + μ + 1$ \\
GKR, $μ$ odd & $3 + ((μ−1)/2)² + 3(μ−1)/2$ \\
\bottomrule
\end{tabular}
\end{center}

\subsection{The GKR in full detail}\label{sec:gen-bus-gkr}

Both verifiers run the same loop (\code{gkr.rs:358-430}, \code{py:433-457}). State: the list
\code{values} of three claimed values (initially $(R, R, R_c)$), the point \code{point}
(initially empty), the combiner $λ$, the level \code{layer} (initially $μ$; the root is level
$μ$, the leaves level 0).

\begin{sloppypar}
\paragraph{Radix and parity.} While $\mathit{layer} > 0$: if \code{layer} is odd the step is one
binary level, else two levels at once (\code{gkr.rs:377}, \code{py:442}). \code{layer} is odd
only at the first iteration, when $μ$ is odd (\code{gkr.rs:378}
\code{debug_assert_eq!(round_count, 0, "only the root-most layer may be binary")}).
Specification: \code{05:85}.
\end{sloppypar}

\paragraph{One iteration, radix 4} (\code{layer} even, $k = μ − \mathit{layer}$ rounds).

{\footnotesize\setlength{\tabcolsep}{4pt}
\begin{longtable}{@{}>{\raggedright\arraybackslash}p{0.09\linewidth}>{\raggedright\arraybackslash}p{0.51\linewidth}>{\raggedright\arraybackslash}p{0.21\linewidth}>{\raggedright\arraybackslash}p{0.13\linewidth}@{}}
\caption{One radix-4 iteration of the GKR verifier.}\label{tab:gen-bus-gkr-iteration}\\
\toprule
Step & What happens & Rust & Python \\
\midrule
\endfirsthead
\toprule
Step & What happens & Rust & Python \\
\midrule
\endhead
\bottomrule
\endfoot
1 & $V$: $\mathit{claim} := \mathit{values}[0] + λ·\mathit{values}[1] + λ²·\mathit{values}[2]$.
  The three trees, in the order push, pull, count, are batched by the powers $1, λ, λ²$
  (\code{05:91-94}).
  & \code{gkr.rs:376}, \code{poly_eval} low degree first,
  \code{primitives/src/}\allowbreak\code{multilinear.rs:245-247}
  & \code{py:443} \\
2 & $k$ sumcheck rounds. Round $t$ ($t = 0 … k−1$) uses the coordinate $\mathit{point}[t]$ as its
  equality factor. & & \\
2, message & $P→V$: \textbf{4 elements}, the coefficients $c_1, c_2, c_3, c_4$ of a polynomial $h$
  of degree 4.
  & \code{gkr.rs:319-326}; read at \code{gkr.rs:399-401}
  \code{next_round_poly(5, claim, Some(equality_point))}
  & \code{py:445} with \code{count = 2**step + 1 = 5} \\
2, derivation & $V$: \textbf{derives} $c_0 := \mathit{claim} + \mathit{point}[t]·(c_1 + c_2 + c_3 + c_4)$.
  This is the round identity $(1 + r)·h(0) + r·h(1) = \mathit{claim}$ solved for $c_0$.
  \textbf{There is no round check to fail}: it holds by construction.
  & \code{fs/transcript.rs:297-302} ``\code{c0 + r·(c1 + … + cd) = claim}, and every \code{ci}
  above \code{c0} was read''
  & \code{py:411-413} \\
2, challenge & $V→P$: $χ_t$; $\mathit{claim} := h(χ_t)$.
  & \code{gkr.rs:402-404}
  & \code{py:424-426} \\
3 & $P→V$: \textbf{12 elements}, the four children of each tree at the point, tree by tree.
  & \code{gkr.rs:406-409}
  & \code{py:447} \\
4 & $V$ \textbf{check}: $\mathit{claim} = Σ_s λ^s · ∏_{c<4} \mathit{children}[s][c]$,
  \textbf{with no equality factor}.
  & \code{gkr.rs:410-413} (\code{LayerMismatch})
  & \code{py:448-449} \\
5 & $V→P$: $u_0$ then $u_1$. $\mathit{values}[s] :=$ the bilinear interpolation of the four
  children, $u_0$ on the low child bit and $u_1$ on the high one.
  & \code{gkr.rs:414-415}; \code{gkr.rs:416-422}
  & \code{py:451}; \code{py:452} \\
6 & $V→P$: a new combiner $λ$.
  & \code{gkr.rs:423}
  & \code{py:453} \\
7 & $\mathit{point} := (u_0, u_1, χ_0, …, χ_{k−1})$; $\mathit{layer} := \mathit{layer} − 2$.
  & \code{gkr.rs:424-425}
  & \code{py:454} \\
\end{longtable}
}

$h$ is the \textbf{cofactor}: the honest
$h_t(Y) = Σ_x \eq(\mathit{point}_{>t}, x) · Σ_s λ^s · ∏_{a,b} \tilde{V}^s_{\mathit{layer}−2}(a, b, χ_{<t}, Y, x)$,
without the factor $\eq(\mathit{point}[t], Y)$ and without the factors
$\eq(\mathit{point}[u], χ_u)$ of the earlier rounds. The running claim is the value of the
cofactor, never multiplied by an equality factor (the ``normalized'' form of \code{gkr.rs:4-5}
``Its normalized eq-trick sumcheck has degree four'').

\paragraph{One iteration, binary} ($\mathit{layer} = μ$ odd, no round): $P→V$ 6 elements, two
children per tree; check $\mathit{claim} = Σ_s λ^s · \mathit{left}_s · \mathit{right}_s$; $V→P$
$u$; $\mathit{values}[s] := (1+u)·\mathit{left}_s + u·\mathit{right}_s$; $V→P$ a new $λ$;
$\mathit{point} := (u)$ (\code{gkr.rs:377-395}, \code{py:442-455} with \code{step = 1}).

\paragraph{The last combiner is never used.} Step 6 runs in every iteration, the last included,
where \code{layer} then becomes 0 and the loop ends. So the number of combiners is the number of
layers plus one, and the last squeeze changes the state of the chain before $ξ$ is drawn and
enters no computation (\code{gkr.rs:423} then \code{gkr.rs:429}; \code{py:453} then
\code{py:457}). The probe confirms it: at $μ = 4$ the GKR draws 9 challenges and $ζ$ is draws 6,
7, 4, 5; draw 8 is unused.

\begin{sloppypar}
\paragraph{How $ζ$ is assembled.} $ζ$ is the \code{point} after the last iteration
($\mathit{layer} = 2$, $k = μ − 2$): $ζ = (u_0, u_1, χ_0, …, χ_{μ−3})$. Every coordinate of $ζ$ is
a challenge of the \textbf{last} layer; no challenge of an earlier layer is a coordinate of $ζ$.
In time order the last layer draws $χ_0, …, χ_{μ−3}$ (which become $ζ_2, …, ζ_{μ−1}$), then
$u_0 = ζ_0$, then $u_1 = ζ_1$, then the unused combiner. Coordinate 0 is the lowest bit of the
leaf index: the Rust test \code{radix_four_roundtrip_at_even_and_odd_depths}
(\code{gkr.rs:479-508}) asserts
\code{proved.values[lane] == mle_eval_e(&leaves[lane], &proved.point)} with \code{mle_eval_e}
folding \code{point[0]} over adjacent pairs (\code{gkr.rs:436-447}), and the probe checks the
same against the pinned Python. For $μ = 1$, $ζ = (u)$; for $μ = 0$ there is no layer and $ζ$ is
empty.
\end{sloppypar}

The probe's table (dossier \file{gt-bus.md}, section H.1): $ζ_k$ is the draw of that number,
counting the GKR's draws from 0.

\begin{center}\small
\begin{tabular}{@{}lll@{}}
\toprule
$μ$ & $ζ_0, ζ_1, ζ_2, …$ & draws \\
\midrule
2 & 1, 2 & 4 \\
3 & 4, 5, 3 & 7 \\
4 & 6, 7, 4, 5 & 9 \\
5 & 10, 11, 7, 8, 9 & 13 \\
8 & 22, 23, 16, …, 21 & 25 \\
\bottomrule
\end{tabular}
\end{center}

\paragraph{How the three trees share challenges.} Entirely: one \code{point}, one sequence of
$χ$, one pair $(u_0, u_1)$ per layer; the three are told apart only by the power of $λ$. The
count tree is shallower; it is padded with leaves $1$ to the depth of the bus trees
(\code{leaf.rs:879-881} \code{count_lay.mu = push_lay.mu}; \code{py:540}, \code{py:578-579},
where the selector of a count block is extended by zero bits to the length of $ζ$).

\subsection{The leaf stacks and the decomposition}\label{sec:gen-bus-leaves}

\paragraph{Blocks, in index order} (\code{layout.rs:349-415}; \code{py:537-540},
\code{py:866-876}). Push side: the state boundary $κ = 0$, $(\mathit{sep}_{\mathrm{ST}}, g⁰, g⁰)$;
the memory seed $κ = κ_{\mathrm{mem}}$,
$(\mathit{sep}_{\mathrm{MEM}}, \mathit{idx}, 1, \mathit{mem}_0, \mathit{mem}_1, \mathit{mem}_2)$;
the bytecode seed $κ = κ_{\mathrm{bc}}$, $(\mathit{sep}_{\mathrm{BC}}, \mathit{idx}, 1, P_3, …, P_{10})$;
then for each table in order each of its pushes, $κ = τ_j$. Pull side: the state boundary
$(\mathit{sep}_{\mathrm{ST}}, g^{2^{κ_{\mathrm{bc}}} − 1}, g⁰)$ (\code{layout.rs:335, 355-358});
the memory finalization with \code{cntfin_mem} in place of $1$; the bytecode finalization with
\code{cntfin_bc}; then the tables' pulls. Count side: for each table in order, one block per
count column, a single coordinate (\code{layout.rs:412-414}). \textbf{The three blocks no table
owns come first.}

\begin{sloppypar}
\paragraph{Placement.} Largest first at aligned offsets, ties by index (\code{witness.rs:67-79}
\code{order.sort_by(|&a, &b| kappas[b]….cmp(…).then(a.cmp(&b)))}, called by
\code{leaf.rs:149-156}; \code{py:305-311}). $\mathit{sel}_b = \mathit{offset}_b ≫ κ_b$, its bits
low first (\code{leaf.rs:407-411}; \code{py:299-302}). The depth is
$⌈\log_2 Σ_b 2^{κ_b}⌉$ with \textbf{no floor} (\code{leaf.rs:153}), unlike the witness stack
(\code{witness.rs:97}).
\end{sloppypar}

\paragraph{Leaves.} Row $z$ of block $b$ holds $β + Σ_i \eq(α, i)·c_{b,i}(z)$
(\code{leaf.rs:226-260}; characteristic 2), the weights $\eq(α, \text{bits of } i)$, bit $k$ of
$i$ against $α_k$ (\code{leaf.rs:89-98}). The count side runs at $α = 0$, $β = 0$, so its leaf is
the count cell (\code{leaf.rs:606-622}; \code{py:573}). Padding leaves are $1$
(\code{leaf.rs:222-224}, \code{gkr.rs:61, 114}).

\paragraph{Decomposition} (the last computation of the phase, U7). With
$ζ_{\mathrm{lo}} = ζ_{<κ_b}$, $ζ_{\mathrm{hi}} = ζ_{≥κ_b}$,
$w_b = \eq(\mathit{sel}_b, ζ_{\mathrm{hi}})$:
\begin{itemize}
\item a block no table owns contributes $w_b·(β + Σ_i \eq(α, i)·v_i)$ where $v_i$ is the
  constant, the index column's extension $∏_k (1 + ζ_k(1 + g^{2^k}))$
  (\code{primitives/src/field/mod.rs:105-113}, \code{py:271-278}; \code{06:97-100}), the extension
  of the public bytecode column, or the value the prover sent for a committed column
  (\code{leaf.rs:440-457}). A committed column already valued at the same point is not sent
  again (\code{leaf.rs:466-471}), which is why the boundary evaluations (U6) are 5 elements and
  not 8;
\item a block of table $j$ contributes nothing to the check: it is accumulated into the form of
  $(\mathit{side}, j)$, coefficient $w_b·\eq(α, i)$ on coordinate $i$ and $w_b·β$ on the constant
  (\code{leaf.rs:416-423}, \code{leaf.rs:367-382}; \code{py:582-587});
\item the padding contributes $1 + Σ_b w_b$ over all the blocks of the side
  (\code{leaf.rs:459-460}; \code{py:589}; \code{05:107});
\item $\mathit{rem}_s := \tilde{V}^s_0(ζ) + (\text{public part} + \text{padding})$
  (\code{leaf.rs:921-924} ``What the tables owe this side: DERIVED, never read'';
  \code{py:590}).
\end{itemize}

The Rust evaluates the eight bytecode columns one by one (\code{leaf.rs:453}); the Python
evaluates the stacked bytecode table once at $(ζ_{<κ_{\mathrm{bc}}}, α)$ (\code{py:566}), as
\code{08:70} writes it. The two are equal when the stacked table is zero outside slots 3 to 10
(\cref{tab:gen-bus-disagreements}, item 6).
